\documentclass[12pt, a4paper]{article}
\usepackage[utf8]{inputenc}
\usepackage[russian]{babel}
\usepackage[T2A]{fontenc}
\usepackage{graphicx}
\usepackage{listings}
\usepackage{xcolor}
\usepackage{hyperref}
\usepackage{geometry}
\usepackage{amsmath}
\usepackage{fancyhdr}
\usepackage{tcolorbox}
\usepackage{tikz}
\usepackage{enumitem}

% Настройка полей
\geometry{left=2cm, right=1.5cm, top=2cm, bottom=2cm}

% Настройка цветов для кода
\definecolor{codegreen}{rgb}{0,0.6,0}
\definecolor{codegray}{rgb}{0.5,0.5,0.5}
\definecolor{codepurple}{rgb}{0.58,0,0.82}
\definecolor{backcolour}{rgb}{0.96,0.96,0.94}
\definecolor{navyblue}{rgb}{0.0, 0.0, 0.5}

% Настройка стиля кода Verilog
\lstdefinestyle{verilogstyle}{
    backgroundcolor=\color{backcolour},   
    commentstyle=\color{codegreen},
    keywordstyle=\color{navyblue}\bfseries,
    numberstyle=\tiny\color{codegray},
    stringstyle=\color{codepurple},
    basicstyle=\ttfamily\footnotesize,
    breakatwhitespace=false,         
    breaklines=true,                 
    captionpos=b,                    
    keepspaces=true,                 
    numbers=left,                    
    numbersep=5pt,                  
    showspaces=false,                
    showstringspaces=false,
    showtabs=false,                  
    tabsize=4,
    frame=single,
    rulecolor=\color{black!20},
    language=Verilog
}

\lstset{style=verilogstyle}

% Красивые блоки для заметок
\newtcolorbox{info}[1]{colback=blue!5!white,colframe=blue!75!black,title=#1}
\newtcolorbox{warning}[1]{colback=red!5!white,colframe=red!75!black,title=#1}
\newtcolorbox{taskbox}[1]{colback=green!5!white,colframe=green!75!black,title=#1}

% Колонтитулы
\pagestyle{fancy}
\fancyhf{}
\rhead{Основы ПЛИС: Sipeed Tang Nano 20K}
\lhead{\leftmark}
\rfoot{Стр. \thepage}

\title{\textbf{Цифровая вселенная на ладони}\\ \large Учебный курс по проектированию цифровых устройств на ПЛИС (FPGA) для старшеклассников}
\author{Лаборатория Цифровой Электроники}
\date{\today}

\begin{document}

\maketitle
\tableofcontents
\newpage

\section{Предисловие для учителя и ученика}
Этот курс создан не просто для того, чтобы научить вас мигать светодиодом. Наша цель — изменить ваше мышление. В школе, изучая информатику, вы привыкли думать \textit{последовательно}: "сделай шаг А, затем шаг Б". 

Но мир вокруг нас — \textit{параллелен}. Пока вы читаете этот текст, ваше сердце бьется, легкие дышат, а за окном едут машины. Все это происходит одновременно. ПЛИС (FPGA) — это технология, которая позволяет создавать цифровые миры, живущие по законам параллельной вселенной.

Здесь мы не пишем программы. Мы проектируем схемы. Добро пожаловать в мир Hardware Engineering.

\newpage

\section{Часть 1. Фундамент цифрового мира}

\subsection{1.1. Аналоговый vs Цифровой сигнал}
Мир вокруг нас аналоговый. Звук — это волна, свет — это спектр, температура меняется плавно. Но компьютеры не понимают "плавно". Они понимают "Да" или "Нет".

\begin{info}{Ключевое понятие: Бит}
\textbf{Бит (bit)} — это минимальная единица информации. Он может быть равен либо \textbf{0}, либо \textbf{1}.
В физическом мире электроники:
\begin{itemize}
    \item \textbf{0 (Логический ноль)} = Низкое напряжение (GND, 0 Вольт). "Лампочка выключена".
    \item \textbf{1 (Логическая единица)} = Высокое напряжение (VCC, 3.3 Вольта). "Лампочка включена".
\end{itemize}
\end{info}

\subsection{1.2. Логические вентили — кирпичики процессора}
Любой, даже самый мощный процессор Intel или Apple M1, состоит из миллиардов простейших элементов — логических вентилей. В ПЛИС мы будем использовать их напрямую.

\subsubsection{NOT (НЕ / Инвертор)}
Самый простой элемент. Он делает все наоборот.
\begin{itemize}
    \item Вход: 0 $\rightarrow$ Выход: 1
    \item Вход: 1 $\rightarrow$ Выход: 0
\end{itemize}
\textit{В Verilog:} \texttt{assign y = \~{}a;}

\subsubsection{AND (И / Конъюнкция)}
Выдает 1, только если \textbf{ВСЕ} входы равны 1. Это как строгий охранник: пустит только если у тебя есть билет И паспорт.
\begin{itemize}
    \item 0 AND 0 = 0
    \item 0 AND 1 = 0
    \item 1 AND 0 = 0
    \item 1 AND 1 = 1
\end{itemize}
\textit{В Verilog:} \texttt{assign y = a \& b;}

\subsubsection{OR (ИЛИ / Дизъюнкция)}
Выдает 1, если \textbf{ХОТЯ БЫ ОДИН} вход равен 1. Это добрый охранник: пустит, если есть билет ИЛИ если ты друг начальника.
\begin{itemize}
    \item 0 OR 0 = 0
    \item 0 OR 1 = 1
    \item 1 OR 0 = 1
    \item 1 OR 1 = 1
\end{itemize}
\textit{В Verilog:} \texttt{assign y = a | b;}

\subsubsection{XOR (Исключающее ИЛИ)}
Очень важный элемент для математики. Выдает 1, если входы \textbf{РАЗНЫЕ}.
\begin{itemize}
    \item 0 XOR 0 = 0
    \item 0 XOR 1 = 1
    \item 1 XOR 0 = 1
    \item 1 XOR 1 = 0
\end{itemize}
\textit{В Verilog:} \texttt{assign y = a \^{} b;}

\newpage

\section{Часть 2. Что такое ПЛИС (FPGA)?}

\subsection{2.1. Расшифровка названия}
\textbf{ПЛИС} — Программируемая Логическая Интегральная Схема.
\textbf{FPGA} — Field-Programmable Gate Array (Матрица вентилей, программируемая пользователем).

Слово "Field" (Поле) здесь означает "на месте эксплуатации". То есть, в отличие от заводских микросхем, структуру которых создают на заводе раз и навсегда, FPGA можно перестраивать прямо у вас на столе бесконечное количество раз.

\subsection{2.2. Архитектура: Что внутри черного квадрата?}
Если посмотреть на кристалл FPGA под микроскопом, мы увидим не процессор, а огромное поле одинаковых ячеек.

\subsubsection{1. LUT (Look-Up Table) — Таблица поиска}
Это сердце ПЛИС. Представьте, что вы не умеете считать $2+2$, но у вас есть шпаргалка, где написано "если спросят 2 и 2, отвечай 4". 
LUT — это маленькая память. Если мы хотим сделать элемент AND, мы просто записываем в эту память таблицу истинности AND. Хотим XOR? Переписываем таблицу.
\textbf{Вывод:} ПЛИС не "вычисляет" логику, она "вспоминает" ответ.

\subsubsection{2. D-Flip-Flop (Триггер) — Элемент памяти}
LUT позволяет вычислять, но не позволяет запоминать. Чтобы сохранить состояние (например, счетчик секунд), нужны триггеры.
Триггер запоминает значение на входе в момент удара "тактового сигнала" и хранит его до следующего удара.

\subsubsection{3. Routing Matrix (Матрица связей)}
Это "провода", которые соединяют LUT и триггеры. При прошивке ПЛИС мы физически переключаем электронные ключи, прокладывая дорожки сигнала.

\subsection{2.3. ПЛИС против Микроконтроллера (CPU)}
\begin{table}[h]
\centering
\begin{tabular}{|l|l|l|}
\hline
\textbf{Характеристика} & \textbf{Микроконтроллер (Arduino/STM32)} & \textbf{ПЛИС (FPGA)} \\ \hline
\textbf{Принцип работы} & Последовательное выполнение команд & Параллельная работа схемы \\ \hline
\textbf{Скорость} & Ограничена тактами на инструкцию & Ограничена скоростью света \\ \hline
\textbf{Сложность} & Легко начать (C++) & Сложнее (Verilog/VHDL) \\ \hline
\textbf{Гибкость} & Жесткая архитектура & Любая архитектура \\ \hline
\textbf{Применение} & Управление меню, простые задачи & Видеообработка, AI, радары \\ \hline
\end{tabular}
\end{table}

\textbf{Аналогия:}
\begin{itemize}
    \item \textbf{CPU} — это один сверхбыстрый повар на кухне. Он режет лук, потом варит суп, потом жарит мясо. Если заказов много, очередь растет.
    \item \textbf{FPGA} — это завод-кухня с конвейером. Один робот только режет лук, другой только варит. Все происходит одновременно. Производительность колоссальная.
\end{itemize}

\newpage

\section{Часть 3. Язык Verilog и философия RTL}

Мы будем использовать **Verilog**. Но прежде чем писать код, нужно понять, что мы вообще делаем.

\subsection{3.1. Это не программирование!}
Verilog — это **HDL** (Hardware Description Language).
\begin{itemize}
    \item \textbf{C++ / Python (Software):} Вы пишете \textit{инструкцию} для процессора. "Возьми число, сложи, положи обратно". Процессор выполняет это шаг за шагом во времени.
    \item \textbf{Verilog (Hardware):} Вы описываете \textit{чертеж схемы}. "Соедини этот провод с входом сумматора, а выход сумматора с триггером". Эта схема существует в пространстве и работает всегда и одновременно.
\end{itemize}

\begin{warning}{Аналогия}
Программа — это рецепт пирога (шаг 1, шаг 2...).
HDL-код — это схема электрической проводки в доме. Свет в спальне и холодильник на кухне работают одновременно, они не ждут друг друга.
\end{warning}

\subsection{3.2. Что такое RTL (Register Transfer Level)?}
Это уровень абстракции, на котором мы работаем. **Уровень Регистровых Передач**.
Вся цифровая электроника сводится к двум вещам:
\begin{enumerate}
    \item **Регистры (Registers):** Это "бассейны", где хранится вода (данные). Они обновляются только по тактовому сигналу (как шлюзы).
    \item **Логика (Combinational Logic):** Это "трубы и фильтры" между бассейнами. Вода течет по ним мгновенно (на самом деле нет, но очень быстро).
\end{enumerate}

**Суть RTL-дизайна:** Мы описываем, как данные перетекают из одного регистра в другой через облака логики (сложение, сравнение, AND/OR) в каждом такте.

\subsection{3.3. От кода к "железу": Синтез}
Когда вы нажимаете кнопку "Run" в среде разработки (IDE), происходит не компиляция, а **Синтез**.
\begin{enumerate}
    \item **Анализ:** Компьютер читает ваш текст.
    \item **Синтез:** Он превращает текст в список логических вентилей (Netlist). "Ага, здесь он написал \texttt{a + b}, значит мне нужно поставить сумматор из 100 транзисторов".
    \item **Place \& Route (Размещение и Трассировка):** Самый сложный этап. Компьютер берет этот список вентилей и пытается физически разложить их на кристалле нашей конкретной ПЛИС (Tang Nano 20K) и соединить проводами. Это как разводка печатной платы, только внутри чипа.
    \item **Bitstream:** Генерация файла прошивки (карты соединений), которая загружается в ПЛИС.
\end{enumerate}

\subsection{3.4. Синтаксис Verilog}

\subsubsection{Модуль — это микросхема}
Любой проект состоит из модулей. Это черные ящики с входами и выходами.

\begin{lstlisting}
module my_chip(
    input  wire clk,   // Вход тактовой частоты
    input  wire btn,   // Вход кнопки
    output reg  led    // Выход на светодиод (регистровый)
);

    // Внутренности схемы (тело модуля)

endmodule
\end{lstlisting}

\subsubsection{Типы данных: wire и reg}
Это самый сложный момент для новичков, но он критически важен.
\begin{itemize}
    \item \textbf{wire (провод):} Просто кусок металла. Он не может хранить информацию. Если вы подали на один конец провода 1, на другом сразу появится 1. Если убрали — пропадет. Используется для соединения блоков.
    \item \textbf{reg (регистр):} Элемент памяти (триггер). Он запоминает значение и хранит его до следующего изменения. В физическом мире это D-триггер.
\end{itemize}

\subsubsection{Присваивания: = против <=}
В Verilog есть два мира:
\begin{itemize}
    \item **Комбинационная логика (\texttt{assign}):**
    \texttt{assign y = a \& b;}
    Работает мгновенно. Как только изменилось \texttt{a}, сразу меняется \texttt{y}. Описывает провода и вентили.
    
    \item **Секвенциальная логика (\texttt{always @(posedge clk)}):**
    Описывает то, что происходит по фронту тактового сигнала (тик часов).
    Здесь мы используем **неблокирующее присваивание (\texttt{<=})**.
    
    \begin{lstlisting}
    always @(posedge clk) begin
        a <= b;
        b <= a;
    end
    \end{lstlisting}
    \textit{Магия:} Значения \texttt{a} и \texttt{b} поменяются местами! Потому что \texttt{<=} означает "запланировать обновление на конец такта". Все обновления произойдут одновременно. В обычном программировании (\texttt{a=b; b=a;}) вы бы потеряли значение \texttt{a}.
\end{itemize}

\newpage

\section{Часть 4. Практикум: Sipeed Tang Nano 20K}

В наших руках плата на базе чипа GW2AR-18.
\begin{itemize}
    \item \textbf{Логических ячеек:} 20 736 (это много, хватит на свой процессор!).
    \item \textbf{Память:} Встроенная SDRAM.
    \item \textbf{Периферия:} HDMI, светодиоды, кнопки.
\end{itemize}

\subsection{Схема платы (Pinout)}
Для выполнения заданий нам нужно знать, куда подключены компоненты:
\begin{itemize}
    \item \textbf{CLK (Тактовый генератор 27 МГц):} Pin 4.
    \item \textbf{LEDs (Светодиоды):} Pins 15, 16, 17, 18, 19, 20.
    \item \textbf{Кнопки:} Pin 87, Pin 88.
\end{itemize}
\textit{Важно:} Светодиоды на этой плате управляются инверсно. Подача "0" включает их, подача "1" выключает.

\newpage
\section{Практические задания}

\subsection{Задание 1: Hello World (Включи свет)}
\begin{taskbox}{Цель}
Научиться создавать проект, описывать простейшую связь "выход-значение" и прошивать ПЛИС.
\end{taskbox}

\textbf{Теория:}
Самая простая схема — это батарейка, подключенная к лампочке. В цифровом мире это константа. Мы хотим, чтобы на ножке, к которой подключен светодиод, всегда был логический ноль (так как схема включения инверсная).

\textbf{Код решения:}
\begin{lstlisting}
module top(
    output led // Объявляем выход
);
    // assign - это "непрерывное присваивание".
    // Представьте, что мы припаяли провод к земле (0).
    // 1'b0 означает: 1 бит, binary (двоичное), значение 0.
    assign led = 1'b0; 
endmodule
\end{lstlisting}

\textbf{Файл ограничений (.cst):}
Здесь мы говорим компилятору: "Сигнал \texttt{led} из кода подключи к физической ножке 15 микросхемы".
\begin{lstlisting}
IO_LOC "led" 15;
IO_PORT "led" IO_TYPE=LVCMOS33; // Уровень логики 3.3 Вольта
\end{lstlisting}

---

\subsection{Задание 2: Мигалка (Blinky)}
\begin{taskbox}{Цель}
Понять концепцию времени и тактовой частоты. Изучить счетчики.
\end{taskbox}

\textbf{Теория:}
У ПЛИС нет понятия "секунда". У неё есть только "такт". Кристалл на плате тикает 27 000 000 раз в секунду (27 МГц).
Чтобы мигнуть раз в секунду, нам нужно досчитать до 13.5 миллионов (полсекунды горит, полсекунды нет).
Для этого нужен \textbf{счетчик} — регистр, который увеличивается на 1 с каждым тактом.

Сколько бит нужно для числа 27 000 000? $2^{24} \approx 16.7$ млн, $2^{25} \approx 33.5$ млн. Значит, нам нужен регистр шириной 25 бит.

\textbf{Код (фрагмент):}
\begin{lstlisting}
reg [24:0] counter; // 25-битный регистр

always @(posedge clk) begin
    counter <= counter + 1; // Увеличиваем на каждом фронте
end

assign led = counter[24]; // Берем самый старший бит
\end{lstlisting}

\textit{Вопрос на засыпку:} Почему мы берем старший бит? Старший бит меняется реже всего. Младший (0-й) меняется 27 млн раз в секунду (мы этого не увидим), 1-й — 13.5 млн раз... 24-й бит будет меняться с частотой примерно $27MHz / 2^{24} \approx 1.6$ Гц.

---

\subsection{Задание 3: Быстро и медленно}
\begin{taskbox}{Цель}
Работа с битами одного счетчика. Делители частоты.
\end{taskbox}

\textbf{Задача:} Подключите два светодиода. Пусть один мигает быстро, а другой медленно.
\textbf{Решение:} Используйте один большой счетчик. Подключите \texttt{led\_fast} к 22-му биту, а \texttt{led\_slow} к 24-му биту счетчика. Это наглядно покажет, как деление частоты на 2 происходит с каждым следующим битом.

---

\subsection{Задание 4: Бинарный счетчик}
\begin{taskbox}{Цель}
Визуализация двоичной системы счисления.
\end{taskbox}

\textbf{Теория:}
Мы привыкли считать в десятичной системе (0-9). Компьютеры считают в двоичной (0-1).
Если вывести состояние счетчика на ряд светодиодов, мы увидим, как "бегут" биты.
\begin{itemize}
    \item 000 (0)
    \item 001 (1)
    \item 010 (2)
    \item 011 (3)
    \item 100 (4)
\end{itemize}

\textbf{Практика:} Выведите старшие 6 бит вашего счетчика на 6 светодиодов платы. Не забудьте про инверсию (\texttt{\~counter}).

---

\subsection{Задание 5: Бегущий огонь (Сдвиговый регистр)}
\begin{taskbox}{Цель}
Изучение операторов сдвига и битовых манипуляций.
\end{taskbox}

\textbf{Теория:}
Сдвиговый регистр — это основа последовательной передачи данных (как USB или Интернет). Представьте шеренгу солдат (биты). По команде каждый делает шаг влево.
В Verilog это оператор \texttt{<<}.
\texttt{regs <= \{regs[4:0], regs[5]\};} — это циклический сдвиг. Мы берем младшие 5 бит и приклеиваем к ним старший бит справа.

---

\subsection{Задание 6 и 7: Усложнение движения}
Реализуйте движение вправо (\texttt{>>}).
А затем "Пинг-понг": добавьте регистр-флаг \texttt{dir} (направление). Если огонек дошел до края, меняйте \texttt{dir} на противоположный. Это простейший пример автомата состояний.

---

\subsection{Задание 8: Кнопка и Ввод}
\begin{taskbox}{Цель}
Обработка внешних сигналов.
\end{taskbox}

\textbf{Теория:}
До этого мы только выводили информацию. Теперь будем читать.
\texttt{input key} — это вход.
Простейшая логика: \texttt{assign led = key;} (или \texttt{\~key} в зависимости от схемы).
Светодиод будет гореть, пока вы держите кнопку.

---

\subsection{Задание 9: Кнопка с фиксацией и Дребезг контактов}
\begin{warning}{Проблема реального мира: Дребезг (Bouncing)}
Механическая кнопка — это пружинка и пластина. Когда вы нажимаете её, контакт не замыкается идеально сразу. Он "дребезжит", замыкаясь и размыкаясь десятки раз за миллисекунды.
Для ПЛИС, работающей на 27 МГц, это выглядит как 100 нажатий подряд!
\end{warning}

\textbf{Задача:} Сделать так, чтобы одно нажатие переключало светодиод (вкл/выкл).
\textbf{Решение:} Нужно написать модуль "Дебаунсер" (Debouncer). Логика: "Если сигнал изменился и остается стабильным в течение 10 мс, то считаем, что нажатие произошло". Только после этого меняем состояние светодиода.

---

\subsection{Задание 10: Сдвиговый регистр с кнопкой}
Объединяем знания. Используйте кнопку (с подавлением дребезга!) как тактовый сигнал для сдвига.
Нажал — огонек сдвинулся. Нажал — сдвинулся еще.

---

\subsection{Задание 11 и 12: ШИМ (PWM) — Имитация аналога}
\begin{taskbox}{Цель}
Как управлять яркостью, если есть только 0 и 1?
\end{taskbox}

\textbf{Теория:}
Широтно-Импульсная Модуляция (PWM). Если включать и выключать светодиод очень быстро (быстрее, чем видит глаз, >100 Гц), то глаз усреднит яркость.
\begin{itemize}
    \item 10\% времени горит, 90\% нет = Тускло.
    \item 90\% времени горит, 10\% нет = Ярко.
\end{itemize}

\textbf{Реализация:} Нужен быстрый счетчик (например, до 255) и регистр сравнения.
\texttt{led = (counter < duty\_cycle) ? 1'b0 : 1'b1;}
В задании 12 (Дышащий светодиод) мы будем плавно менять \texttt{duty\_cycle} туда-сюда.

---

\subsection{Задание 13: Светофор (Конечный автомат / FSM)}
\begin{taskbox}{Цель}
Основа всей сложной логики. Finite State Machine.
\end{taskbox}

\textbf{Теория:}
Любое сложное устройство можно описать как набор состояний и переходов между ними.
\begin{itemize}
    \item Состояние 1: RED. Горит красный. Ждем 3 сек $\rightarrow$ переход в RED\_YELLOW.
    \item Состояние 2: RED\_YELLOW. Горят оба. Ждем 1 сек $\rightarrow$ переход в GREEN.
    \item Состояние 3: GREEN. Горит зеленый. Ждем 3 сек $\rightarrow$ переход в YELLOW.
    \item ...
\end{itemize}
В Verilog это делается через конструкцию \texttt{case (state)}. Это самый профессиональный способ писать код.

---

\subsection{Задание 14: Логические вентили на практике}
Используйте две кнопки как входы A и B.
\begin{itemize}
    \item LED1 = A AND B (Загорится, только если зажать обе).
    \item LED2 = A OR B (Загорится, если нажать любую).
    \item LED3 = A XOR B (Загорится, если нажата только одна из двух).
\end{itemize}
Это отличная демонстрация булевой алгебры в железе.

---

\subsection{Задание 15: Игра на реакцию}
\textbf{Финальный проект.}
Огонек бегает влево-вправо с большой скоростью.
Цель игрока: нажать кнопку ровно в тот момент, когда горит центральный светодиод.
\begin{itemize}
    \item Если попал: Огонек останавливается и радостно мигает (Победа).
    \item Если промазал: Огонек останавливается и горит красным (Поражение).
    \item Сброс игры по второй кнопке.
\end{itemize}
Здесь вам понадобятся: FSM (состояния игры: IDLE, RUN, WIN, LOSE), счетчики, дебаунсинг кнопок и логика сравнения.

\section{Заключение}
Пройдя эти 15 заданий, вы перестали быть просто пользователями компьютеров. Вы стали их создателями. Вы научились управлять электронами на самом низком уровне абстракции.
Дальше перед вами открыт путь к созданию своих видеокарт, нейропроцессоров и контроллеров роботов. Удачи!

\end{document}
